\Artigo%
{historia}%
{Sobre como Schrödinger case morre no concello da Estrada}%
{Manuel Galán Rodríguez}%
{}

\begin{multicols}{2}

Seguro que todos os que estades a ler isto escoitastes falar algunha vez de
Erwin Schrödinger, mesmo antes de terdes pisado un pé na Facultade de Física.
Este destacado físico austríaco é amplamente coñecido polo seu famoso
experimento do Gato de Schrödinger. Menos coñecidas, se cadra, son as súas
falcatruadas amorosas nun hotel en Arosa (Suíza), mais hoxe vou detallar como
foi a súa estada por Santiago de Compostela en 1934.

%O de estada non é unha errata por estadía, en galego normativo dise así.

Desde 1911 véñense realizando as conferencias Solvay, unha serie de congresos
sobre física e química realizados en
Bélxica. Foi nunha delas, en setembro de 1933, onde Blas Cabrera, que terá o
papel de Jorge Mira na nosa historia, tivo o pracer de ser o primeiro físico
español en participar, e tamén se cre que coñeceu alí en persoa a Schrödinger.
De alí a uns meses, Cabrera foi nomeado reitor da Universidade Internacional de
Verán de Santander, co que convidou o Nobel a impartir un curso nela en agosto
de 1934.

\begin{figure}[H]
    \centering
    \includegraphics[width = \linewidth]{revistas/005/imaxes/schrodinger.jpg}
    \caption{
        Erwin Schrödinger dando a conferencia o día 3 de agosto de 1933 no
        Paraninfo (Santiago de Compostela). (Fonte: GCiencia)
    }
    \label{charla}
\end{figure}

Casualmente, uns poucos días antes, tiña lugar en Galiza o XIV Congreso da Asociación
Española para el Progreso de las Ciencias\footnote{Concretamente, as
conferencias estaba previsto que estivesen repartidas entre A Coruña, Santiago
de Compostela, Pontevedra e Vigo.}, que reuniría máis de trescentos
congresistas durante unha semana, comezando o mércores 1 de
agosto de 1933. Por tanto, Cabrera tomou a decisión de convidar tamén a
Schrödinger\footnote{Crese que foi quen o invitou, pero non se sabe con
completa seguridade.}, para que dese alí unha charla sobre mecánica cuántica
ondulatoria (Figura~\ref{charla}) o venres 3 de agosto ás 18:00 no Paraninfo
(Santiago de Compostela).

O discurso tratou da revolución que supuxo a mecánica
cuántica na física, onde destacou o papel
das medicións na mecánica cuántica, e as súas diferenzas coa mecánica clásica.
Máis tarde, tratou sobre as nocións de espazo e tempo dunha partícula, onde
critica o concepto do sólido ríxido ideal: na mecánica cuántica só sistemas
infinitamente pesados poden considerarse ríxidos, pois os niveis enerxéticos
son finitos. A conferencia remata resaltando a dificultade para reconciliar as
dúas teorías da física moderna: a relatividade e a mecánica cuántica.

O devandito congreso contou coa presenza
de xente de todas as disciplinas\footnote{De feito, físicos eran máis ben
poucos, pois Schrödinger comezou a súa intervención desculpándose do contido da
súa comunicación, moi avanzado para a audiencia, poñendo como escusa a pouca
antelación coa que o convidara Cabrera.}, como destacados membros do Seminario
de Estudos Galegos, tales como Filgueira Valverde, Cuevillas ou Otero Pedrayo.
Non obstante, contou coa ausencia dalgúns dos principais químicos da época como
Fernando Calvet ou Isidro Parga Pondal, profesores na USC por aquela altura.

Porén, a prensa galega parece que fixo caso case omiso da conferencia do
brillante físico austríaco daquela sexta feira, se cadra por un aviso tardío da
súa presenza. O \textit{La Voz de Galicia} resaltaba na edición do día
seguinte, do sábado 4, o protagonismo de Ortega y Gasset na segunda xornada,
desenvolta desta vez na Coruña\footnote{Espóiler: ao final a Ortega y Gasset non
lle deu tempo a chegar.}. A mención a Schrödinger só se reducía a «el profesor
francés M. Sehroedinfer (sic)»\footnote{Repárese que se equivocaron tanto no
nome como no xentilicio.}. Este mesmo xornal, o domingo, incorporaría unha
pequena fotografía da conferencia, errando de novo no complicado nome do
protagonista da nosa historia.

A contraportada do \textit{El Pueblo Gallego} dese sábado describía en detalle
o congreso, mais deixando a Schrödinger nunha posición inferior:

\begin{adjustwidth}{1cm}{0.6cm}
    \textit{SECCIÓN DE FÍSICA Y QUÍMICA Disertó ayer en esta sección el insigne premio
    Nobel, profesor Schrodinger, a quien sus ecuaciones diferenciales
    interpretativas de las teorías mecánico ondulatorias, así como sus matrices
    curicas (sic) harían suficientemente famoso, si otros muchos trabajos no
    hubieran conseguido hacerlo destacar. A su simpatía personal une el joven
    % \footnote{Iso de \textit{'joven'} relativamente, unha semana despois faría 46 anos.}
    sabio la especial de verse exilado del país que le vio nacer por los
    bárbaros odios de raza que no se detienen ante las cumbres de las ciencias.
    Tan solo en esta crónica daremos brevemente la nota periodística de la
    emoción enorme que produjo la presencia de este sabio en el auditorio que
    ayer le escuchó en la Sección tercera.}
\end{adjustwidth}

O sábado día 4, as conferencias do congreso trasladaríanse até A Coruña, onde
Cabrera realizaría unha intervención. Non sabemos se Schrödinger acompañou a
Cabrera até alí, pois a prensa da época non ofrece a suficiente información.
Posibelmente si, pois ese mesmo día a mediodía estaba prevista unha
intervención de Ortega y Gasset criticando a «barbarie do especialismo» na
ciencia, e a Schrödinger encantáballe a filosofía.

No entanto, a chegada de Ortega y Gasset atrasaríase un día, o que impediu que
tivese lugar a devandita conferencia. Pola contra, o público si puido gozar
doutras intervencións, como as do propio Cabrera, Otero Pedrayo, as do
historiador Rafael Altamira ou as do vizconde de Eza, presidente da asociación que
organizou o congreso; todas realizadas no teatro Rosalía de Castro. Daquela,
Schrödinger non falaba español, mais si francés, italiano, latín e inglés (e
alemán, evidentemente) que lle terían permitido seguir moitas das conferencias.

A seguinte xornada do congreso era o domingo. Ese día estaban previstas dúas
excursións polo val do Ulla, ás que asistirían, a pesar do mal tempo, unhas 250
persoas. Unha destas saídas era ao pazo de Santa Cruz de Ribadulla (Vedra), e a
outra, que acabou en traxedia, foi ao pazo de Oca (A Estrada), á que asistirían
unhas setenta persoas.

De súpeto, mentres estaban os convidados nun bufete organizado polos marqueses
de Camarasa, caeu o chan do salón (Figura~\ref{accidente}), precipitándose os
alí presentes desde unha altura de 5 m, deixando gravemente feridas máis de
cincuenta persoas e unha morte.
En definitiva, un desafortunado suceso que podería chegar pasar calquera día
na nosa facultade (esperemos que non), tal e como está construída.

Non sabemos con seguridade se Schrödinger quedou ese domingo en Santiago de
Compostela ou se foi até A Estrada ou até Vedra, pois a prensa da época non
achega a suficiente información. O máis probábel que estivese en todos eses sitios
e en ningún ao mesmo tempo. Bromas á parte, posibelmente quedase en Santiago, pois
dese domingo data unha carta escrita en follas con \textit{Hotel Compostela} na
cabeceira, conservada no arquivo Zubiri, onde se pode ver que estaba a preparar
as conferencias que daría máis tarde. O que é seguro é que non formou parte do
medio cento de feridos no accidente.

\begin{figure}[H]
    \centering
    \includegraphics[width = \linewidth]{revistas/005/imaxes/accidente.jpg}
    \caption{
        Imaxe do tráxico suceso no pazo de Oca (A Estrada) acontecido a 5 de
        agosto de 1933. (Fonte: GCiencia)
    }
    \label{accidente}
\end{figure}

Evidentemente, e dada a traxedia, o congreso quedou suspendido. Estaba previsto
que rematase no teatro García Barbón en Vigo, coa intervención de Gregorio
Marañón e de Niceto Alcalá-Zamora, o presidente da República. Polo que a
Schrödinger respecta, este colleu rumbo cara a Santander, onde ofrecería seis
conferencias. Nunca volvería a Galiza, aínda que si volveu a España dous anos
máis tarde, acompañado da súa esposa, nunha viaxe duns oito mil quilómetros a
bordo dun BMW. Cabrera e outros compañeiros tentaron procurar a posibilidade de
outorgarlle unha cátedra na Universidade de Madrid, sen éxito. A crise política
naqueles anos en España e a posterior ditadura franquista provocaron que até o
propio Cabrera tivese que verse exiliado.

Xa para finalizar, que tería pasado se o accidente daquel domingo 5 de agosto
de 1933 tivese sido fatal para Schrödinger? Que consecuencias tería para a
física? O que é seguro é que o famoso experimento mental do Gato de Schrödinger
non podería ter chegado a idearse, pois este desenvolveuse en 1935, co que
xamais podería acabar sendo substituído polo paradoxo do Pazo de Oca.

\nocite{galan1}
\nocite{galan2}
\nocite{galan3}

\printbibliography

\end{multicols}
